% !TEX root = ../main.tex
\section{问题重述}

\subsection{问题背景}

在无线电频谱管理和危险环境排查中，需要根据测向设备获得的方位信息快速缩小干扰源的位置范围。测向机给出的示向度受到当地电磁环境和仪器精度影响，与真实方位角之间存在不超过 $1^\circ$ 的误差；同时，各干扰源的有效接收半径未知，但均位于 $1000$ 至 $1500$ 米之间。因此，定位方法不仅要处理角度误差，还要保证新检测位置能够可靠接收到信号。

\subsection{问题内容}

\textbf{问题一：}根据若干检测点坐标及同一干扰源在这些点处的示向度，构造交会定位多边形，给出该多边形直径的计算算法，并判断以定位区域直径为直径的圆能否覆盖整个定位区域。

\textbf{问题二：}已知某全向干扰源在第一个检测点处的示向度，根据第一次观测形成的位置候选区域，给出兼顾接收可靠性和交会定位精度的第二检测点选择策略，并明确第二检测点的候选区域。

\textbf{问题三：}在半径 $1800$ 米的圆域内，对 $10$--$16$ 个频道、位置及有效接收半径未知的全向干扰源，设计自动搜索、定位与清除策略，使全部干扰源被清除且总时间尽可能短。
